\section{Self-check}

\subsection{Why PCA must be fitted on the training images only}

\subsection{Why a linear projection cannot increase linear separability, and how it can still help}

\subsection{What the first eigen-hand should look like}

\section{Components}

\begin{figure}[H]
  \centering
  \includegraphics[width=\textwidth]{stage2/explained_variance.pdf}
  \caption{Explained variance ratio of the first 50 components (left) and cumulative explained variance up to 500 components (right).}
  \label{fig:s2-evr}
\end{figure}

\begin{figure}[H]
  \centering
  \includegraphics[width=\textwidth]{stage2/eigenhands.pdf}
  \caption{Mean training image and the first ten principal components displayed as $28\times28$ images.}
  \label{fig:s2-eigenhands}
\end{figure}

\begin{figure}[H]
  \centering
  \includegraphics[width=0.8\textwidth]{stage2/reconstructions.pdf}
  \caption{Four validation images (left column) reconstructed from 10, 50 and 150 components.}
  \label{fig:s2-reconstructions}
\end{figure}

\section{Number of components}

\begin{table}[H]
  \centering
  \caption{Accuracy and training time for raw pixels and for different numbers of principal components.}
  \label{tab:s2-k-sweep}
  \input{../tables/s2_k_sweep.tex}
\end{table}

\begin{figure}[H]
  \centering
  \begin{subfigure}[t]{0.48\textwidth}
    \includegraphics[width=\textwidth]{stage2/k_sweep_acc.pdf}
    \caption{Validation accuracy.}
  \end{subfigure}\hfill
  \begin{subfigure}[t]{0.48\textwidth}
    \includegraphics[width=\textwidth]{stage2/k_sweep_time.pdf}
    \caption{Training time.}
  \end{subfigure}
  \caption{Validation accuracy and training time of the softmax classifier as a function of the number of components $k$. Dashed line: raw pixels.}
  \label{fig:s2-k-sweep}
\end{figure}

\section{Noise}

\begin{table}[H]
  \centering
  \caption{Test accuracy on noisy test images for the raw-pixel classifier and the PCA classifier.}
  \label{tab:s2-noise}
  \input{../tables/s2_noise.tex}
\end{table}

\begin{figure}[H]
  \centering
  \begin{subfigure}[t]{0.58\textwidth}
    \includegraphics[width=\textwidth]{stage2/noise_examples.pdf}
    \caption{A test image with increasing noise (top) and its PCA reconstruction (bottom).}
  \end{subfigure}\hfill
  \begin{subfigure}[t]{0.38\textwidth}
    \includegraphics[width=\textwidth]{stage2/noise_acc.pdf}
    \caption{Test accuracy against noise level.}
  \end{subfigure}
  \caption{Effect of Gaussian noise on the test images.}
  \label{fig:s2-noise}
\end{figure}

\section{Prediction vs.\ result}

\subsection{Question 1: what do the components look like?}

\subsection{Question 2: how many components does the classifier need?}

\subsection{Question 3: does PCA really remove noise?}

\section{AI-use note}
